\section{Results}\label{sec:results}

\subsection{Do the named sliders do what their names say?}

First the yardstick. Content enjoyment averages \fmaCE{} over \nReal{} real recordings, \aceCEzero{} for unsteered ACE-Step clips and \mainCEzero{} for Stable Audio Open. A slider that drops a clip well below the real-music mean has left the region where the change it made is worth having.

Table~\ref{tab:ace} and Figure~\ref{fig:ace} report eight ACE-Step prompt-pair sliders over nine positions. Where a descriptor was assigned in advance it rises with the slider, and MuQ-MuLan agrees with CLAP on the direction for all eight. The tag panel says the same in words: the mood slider raises \emph{happy} and lowers \emph{dark-toned}. Enjoyment stays above the real-music mean across the range for every slider but tension, which collapses past $+1$, and that is where its usable span ends. SongEval, trained on different data, tells the same story: unsteered clips score 2.77 in musicality against 2.78 for real recordings, the usable span it implies is identical for six of the eight sliders and half a step wider for the other two, and the two predictors agree over clips with Spearman 0.52--0.67. The recipe carries over unchanged to Stable Audio 3 small, where six sliders at 600 iterations all move their descriptor ($\rho$ 0.34--0.84, four with at least 96\% of trajectories ordered), with the ungated quality cost of a 50-step model.

\begin{table*}[t]\centering\small
\begin{tabular}{llrrrrrrrr}
\toprule
Slider & Descriptor & $\rho$ & Ordered & MuQ $\rho$ & Span & Moved & Kept & CE$_0$ & CE$_\text{ends}$ \\
\midrule
brightness & centroid\_oct & 0.97 & 1.00 & 0.80 & -1.0 to +2.0 & 4.41 & 0.74 & 6.95 & 6.04 \\
density & onset\_rate & 0.78 & 0.92 & 0.86 & -0.5 to +2.0 & 1.25 & 0.84 & 6.95 & 6.09 \\
mood & majorness & 0.39 & 0.75 & 0.88 & -1.0 to +2.0 & 0.87 & 0.79 & 6.95 & 6.65 \\
harmony & harmonic\_change & 0.41 & 0.76 & 0.73 & -1.5 to +2.0 & 0.84 & 0.83 & 6.95 & 6.64 \\
groove & -- & -- & -- & 0.56 & -1.5 to +2.0 & -- & 0.79 & 6.95 & 6.88 \\
ensemble & pc & 0.60 & 0.90 & 0.77 & -1.5 to +2.0 & 0.92 & 0.80 & 6.95 & 6.56 \\
tension & -- & -- & -- & 0.86 & -1.5 to +1.0 & -- & 0.83 & 6.95 & 5.64 \\
melody & key\_clarity & 0.48 & 0.78 & 0.89 & -1.0 to +2.0 & 0.97 & 0.80 & 6.95 & 6.47 \\
\bottomrule
\end{tabular}
\caption{Prompt-pair sliders on ACE-Step 1.5 XL turbo. $\rho$: mean rank correlation between position and descriptor. Ordered: share of trajectories with correctly ordered ends. MuQ $\rho$: the same correlation for the MuQ-MuLan direction score. Span: usable span. Moved: descriptor change over the span in standard deviations of unsteered clips. Kept: CLAP similarity to the unsteered clip at the ends of the span. CE: content enjoyment unsteered and at the extreme positions.}\label{tab:ace}
\end{table*}


\begin{figure*}[t]\centering
\includegraphics[width=0.74\linewidth]{ace/response.png}
\caption{ACE-Step: change in each slider's descriptor against position, in standard deviations of unsteered clips, for prompt-pair sliders (blue) and sliders trained between two sets of clips (pink). Shaded bands are 95\% intervals over trajectories.}\label{fig:ace}
\end{figure*}

\subsection{Four ways to move the same attribute}

Table~\ref{tab:methods} runs the alternatives on the same prompts and seeds. The trained slider reproduces prompt-pair guidance, which is after all its own training target, with one forward pass per step instead of three: it scores higher on brightness and harmony and lands inside the confidence interval everywhere else. Activation steering turns out to be a real competitor. It is within 0.05 of the slider on four of six descriptors and keeps more of the piece (0.83--0.87 against 0.74--0.84), while the slider moves its descriptor further on brightness, ensemble, and harmony. Interpolating the text conditioning has the lowest correlation on every attribute and breaks the audio at the ends.

\begin{table}[t]\centering\footnotesize\setlength{\tabcolsep}{2.5pt}
\begin{tabular}{lrrrrrrrr}
\toprule
 & \multicolumn{2}{c}{Slider} & \multicolumn{2}{c}{Guidance} & \multicolumn{2}{c}{Act. steer.} & \multicolumn{2}{c}{Interp.} \\
Slider & $\rho$ & MuQ & $\rho$ & MuQ & $\rho$ & MuQ & $\rho$ & MuQ \\
\midrule
brightness & 0.97 & 0.80 & 0.89 & 0.81 & 0.93 & 0.81 & 0.49 & 0.35 \\
density & 0.78 & 0.86 & 0.74 & 0.81 & 0.71 & 0.86 & 0.10 & 0.38 \\
mood & 0.39 & 0.88 & 0.45 & 0.87 & 0.41 & 0.88 & 0.25 & 0.63 \\
harmony & 0.41 & 0.73 & 0.08 & 0.68 & 0.24 & 0.72 & -0.01 & 0.23 \\
groove & -- & 0.56 & -- & 0.71 & -- & 0.58 & -- & 0.12 \\
ensemble & 0.60 & 0.77 & 0.55 & 0.75 & 0.58 & 0.67 & 0.25 & 0.25 \\
tension & -- & 0.86 & -- & 0.88 & -- & 0.91 & -- & 0.69 \\
melody & 0.48 & 0.89 & 0.51 & 0.83 & 0.53 & 0.83 & 0.31 & 0.38 \\
\midrule
CE at ends & \multicolumn{2}{c}{6.37} & \multicolumn{2}{c}{6.54} & \multicolumn{2}{c}{6.52} & \multicolumn{2}{c}{4.95} \\
Passes per step & \multicolumn{2}{c}{1} & \multicolumn{2}{c}{3} & \multicolumn{2}{c}{1} & \multicolumn{2}{c}{1} \\
\bottomrule
\end{tabular}
\caption{Four ways to move the same attribute on ACE-Step, same prompts and seeds. $\rho$: rank correlation of position with the waveform descriptor (-- where none was assigned); MuQ: with the MuQ-MuLan direction. CE at ends: mean content enjoyment at the two extreme positions, averaged over sliders.}\label{tab:methods}
\end{table}


\subsection{The embedding says yes, the waveform says no}

On Stable Audio Open, sliders trained from the prompt pairs for note density and for percussion move the output along their CLAP direction at every setting. Neither moves its descriptor: the rank correlation with onset rate or percussive share never exceeds \weakRho{}. Scored by embedding alone, both would have been reported as working. The descriptor has to be sound as well. With a tempo estimate prone to octave errors the tempo slider scored $\rho=-0.01$; with a beat tracker it scores $0.57$ and moves the mean tempo from 90 to 167 BPM.

Something similar shows up when the training recipe changes. Retraining the mood slider with rank 16, with a doubled guidance multiplier, on cross-attention only, or for three times as long leaves the CLAP direction score at 0.92--0.93 in every case, while the descriptor correlation ranges from 0.29 to 0.52. The embedding cannot tell these sliders apart; the waveform can.

\subsection{Leaving the first steps alone}

On the 50-step model a slider at $\pm1$ takes the piece with it. Switched on only from step \gateSteps{}, the brightness slider keeps \gateRangeOn{} of its \gateRangeOff{} standard deviations of range, similarity to the unsteered clip rises from \gateKeepOff{} to \gateKeepOn{}, and the loss in enjoyment falls from \gateLossOff{} to \gateLossOn{} points (Table~\ref{tab:gating}). The earliest steps decide the layout of the piece; a slider that leaves them alone keeps it.

\begin{table}[t]\centering\small\setlength{\tabcolsep}{4pt}
\begin{tabular}{rrrrr}
\toprule
On after & Moved & CLAP & Chroma & CE \\
\midrule
0 & 2.37 & 0.67 & 0.68 & 5.15 \\
10 & 2.11 & 0.71 & 0.71 & 5.34 \\
14 & 1.98 & 0.74 & 0.74 & 5.45 \\
21 & 1.41 & 0.85 & 0.84 & 5.87 \\
29 & 0.42 & 0.97 & 0.97 & 6.08 \\
\bottomrule
\end{tabular}
\caption{Timestep gating for the brightness slider on Stable Audio Open at positions $\pm1$: descriptor moved (std), CLAP and chroma similarity to the unsteered clip, content enjoyment. Unsteered CE is 6.16.}\label{tab:gating}
\end{table}


\subsection{What else moves, and cancelling it}

Sliders are not clean. The mood and melody sliders both raise the spectral centroid by 0.8 standard deviations per unit, and several raise loudness. A measured leak can be cancelled with a second slider: adding the brightness slider at $-0.53$ and $-0.55$ times their position brings the leak to $-0.02$ and $+0.09$ and raises similarity to the unsteered clip at $\pm1$ from 0.85 to 0.90 and from 0.87 to 0.92. The CLAP direction correlation drops too, from 0.92 to 0.69 and from 0.93 to 0.72, which suggests that part of what the embedding scores as happier or more melodic was simply brightness.

\subsection{Axes that real music chose}

Fitted on two random halves of the recordings and matched one to one, the leading eight axes agree with mean cosine 0.98 (MuQ-MuLan, principal), 0.94 (CLAP, principal), 0.90 (MuQ-MuLan, independent), and 0.83 (CLAP, independent). Of 1{,}024 sparse-autoencoder features on CLAP, only 9 are found by two seeds. A harder test refits the axes on 47{,}942 further recordings from a different slice of the archive: the leading subspace carries over (the top 32 directions overlap 97\%) while single axes rotate inside it, and the matched cosine of the MuQ-MuLan independent components falls to 0.61.

The independent components of MuQ-MuLan are the ones that read like music. One runs from quiet, minor-key, and dreamy to aggressive, dry, and rhythmic; another from dark and distorted to playful and simple. These are recognisably the arousal and valence of the circumplex model \citep{russell1980circumplex}, and nobody asked for them. Others contrast jazz with electronic music, plucked and bowed strings with synthesizer and choir, and piano with organ and rock. Arousal, valence, and the jazz--electronic axis reappear in the second corpus with the same tags (cosine 0.65--0.76); the strings--synthesizer axis does not.

\begin{figure*}[t]\centering
\includegraphics[width=0.62\linewidth]{figures/coverage.png}
\caption{How much of each real-music axis \covClips{} ACE-Step clips span, as a fraction of the spread of real recordings, over all 648 prompts (blue) and for one prompt with many seeds (orange). Each axis is labelled by the tags at its two ends.}\label{fig:coverage}
\end{figure*}

\subsection{How much of real music does a generator cover?}

Along every one of these axes the generator is narrower than real music (Figure~\ref{fig:coverage}). Over 648 prompts crossing genre, instrument, and mood, ACE-Step spans \covLo{}--\covHi{}\% of the real spread; seeds of a single prompt span \covSeedLo{}--\covSeedHi{}\% of what recordings of one genre span. The average generated clip sits 1.35 standard deviations toward the quiet end of the arousal axis and 1.36 toward the playful end of the valence axis. On CLAP's first three principal axes ACE-Step covers 35--39\% and Stable Audio Open 48--52\%, from different training data and a different architecture. The top eight principal directions of the generated corpus contain only \covOverlap{}\% of the top eight of the real one: the directions along which generated clips differ most are mostly not the ones along which recordings do.

\IfFileExists{sections/results_axes.tex}{\smallskip\noindent\textbf{Sliders along real-music axes.} Seven axes were turned into sliders with the set trainer, the sets being the top and bottom 20\% of 24 seeds per prompt along the axis (Table~\ref{tab:axes}). All \axN{} follow their axis, with rank correlations between \axRhoLo{} and \axRhoHi{}. The arousal slider moves one prompt and seed 1.19 standard deviations of real music along its axis, more than the standard deviation of the model's whole output over 648 prompts (0.83) and 2.1 times the standard deviation that seeds of one prompt produce. Similarity to the unsteered clip is \axKeepLo{}--\axKeepHi{} at $\pm1$, enjoyment is at or above the unsteered level, and musicality drops by at most 0.14.

\begin{table}[t]\centering\footnotesize\setlength{\tabcolsep}{2.5pt}
\begin{tabular}{llrrrrrrr}
\toprule
Axis & From & $\rho$ & Ord. & Real & Seeds & Kept & CE & Mus. \\
\midrule
acoustic/electr. & sets & 0.55 & 0.93 & 0.64 & 1.48 & 0.86 & 7.11 & 2.68 \\
arousal & sets & 0.82 & 0.99 & 1.19 & 2.14 & 0.84 & 7.13 & 2.70 \\
classical/funk & sets & 0.64 & 0.96 & 0.77 & 1.65 & 0.85 & 7.12 & 2.69 \\
jazz/electronic & sets & 0.61 & 0.94 & 0.61 & 1.23 & 0.87 & 7.09 & 2.70 \\
piano & sets & 0.63 & 0.96 & 0.75 & 1.42 & 0.87 & 7.08 & 2.69 \\
strings/synth & sets & 0.40 & 0.82 & 0.44 & 0.81 & 0.87 & 7.09 & 2.69 \\
valence & sets & 0.50 & 0.93 & 0.62 & 1.23 & 0.87 & 7.07 & 2.73 \\
\bottomrule
\end{tabular}
\caption{Sliders along axes of real music on ACE-Step, positions $-1$ to $+1$. From: a prompt pair made of the axis's tags, or two sets of the model's clips. $\rho$, Ord.: rank correlation between position and the output's projection on the axis, and share of trajectories with $-1$ and $+1$ in the right order. Real, Seeds: movement along the axis between $-1$ and $+1$ in standard deviations of real recordings and of one prompt's seeds. Kept: CLAP similarity to the unsteered clip. CE, Mus.: enjoyment and SongEval musicality at $\pm1$ (unsteered: 6.95 and 2.77).}\label{tab:axes}
\end{table}


\smallskip\noindent\textbf{Each trainer moves what it was trained on.} Sorting clips by a waveform descriptor, residualised against loudness, brightness, and enjoyment, gives sliders that are selective. Over positions $-1$ to $+1$ the set-trained harmony slider reaches $\rho=0.72$ on harmonic change rate with selectivity 6.6, against 0.27 and 0.8 for the prompt-pair slider; density and ensemble show 4.7 and 2.4 against 1.5 and 1.5. The embedding tells the opposite story: the MuQ-MuLan text direction for harmony follows the prompt-pair slider ($\rho=0.63$) and ignores the set-trained one ($-0.07$). Neither kind of score can stand in for the other. Set-trained sliders stop at $\pm1$: beyond it similarity falls below 0.72 and musicality drops by 0.3--0.5 in both directions while CE stays flat, a failure only the second predictor shows. Tempo did not become a slider this way ($\rho=0.06$).
}{}

\subsection{What did not work}

\textbf{Axes inside the generator.} We hoped the generator's own activations would give musical axes without any embedding model. Eight principal axes carry 90\% of the within-prompt variance of ACE-Step's cross-attention outputs and reproduce across disjoint prompt sets (0.82). Steering along the first two raises level and lowers spectral flatness ($\rho=0.77$ and $-0.64$) and moves tags from lo-fi toward complex and live: these are production axes. The other six leave at least 95\% similarity to the unsteered clip at twelve natural standard deviations and act differently from prompt to prompt (consistency 0.10--0.21). Unlike the image case \citep{haas2024directions}, the leading directions here describe polish; arousal, valence, and instrument contrasts only appeared in embeddings of real music.

\textbf{Training between sets of real recordings.} With $H$ and $L$ drawn from the recordings themselves, the slider damaged the clip equally in both directions. A penalty tying the mean of its predictions at $\pm1$ to the frozen model removed the damage and left almost no steering: tempo did not move and major-minor fit went from 0.09 to 0.15. Real recordings sit off the generator's distribution and the mismatch swamps the loss, so real music served to find the axes and the generator's own clips to train along them.
